% =====================================================================
%  CAMPOS CLÁSSICOS — Notas de estudo (início de mestrado em Física Teórica)
%  Transcrição organizada e revisada a partir das notas manuscritas.
%
%  Compilar com:  pdflatex campos_classicos.tex   (rodar 2x)
%  Fonte: Computer Modern (padrão do LaTeX, embutida no PDF).
% =====================================================================

\documentclass[12pt,a4paper]{article}

% ---------------------------------------------------------------------
%  CODIFICAÇÃO E FONTE
%  Computer Modern em T1 (fontes EC) — acentos do português corretos.
%  (Sem babel-português nesta máquina; a hifenização usa o padrão.
%   Para hifenização em pt-BR, instale 'babel-portuguese' e descomente
%   a linha \usepackage[brazil]{babel} mais abaixo.)
% ---------------------------------------------------------------------
\usepackage[T1]{fontenc}
\usepackage[utf8]{inputenc}
% \usepackage[brazil]{babel}   % <- habilite se tiver o pacote instalado

% ---------------------------------------------------------------------
%  MATEMÁTICA
% ---------------------------------------------------------------------
\usepackage{amsmath}
\usepackage{amssymb}
\usepackage{mathtools}

% =====================================================================
%  BLOCO DE FORMATAÇÃO  ------  (altere aqui tudo o que for de layout)
%  Ajustado às regras da ABNT (NBR 14724): A4, margens 3/3/2/2 cm,
%  corpo 12 pt, espaçamento 1,5, recuo de parágrafo 1,25 cm, texto
%  justificado, seções numeradas.
% =====================================================================
\usepackage[a4paper,left=3cm,top=3cm,right=2cm,bottom=2cm]{geometry}
\usepackage{setspace}
\onehalfspacing                        % espaçamento entre linhas: 1,5
\setlength{\parindent}{1.25cm}         % recuo da 1a linha do parágrafo
\usepackage{indentfirst}               % recua também o 1o parágrafo (ABNT)
\usepackage{titlesec}                  % controle simples dos títulos
\titleformat*{\section}{\bfseries\large}       % seção primária: negrito
\titleformat*{\subsection}{\bfseries\normalsize}
\emergencystretch=2em                  % evita "overfull" na justificação

% ---------------------------------------------------------------------
%  MARCAÇÃO DE REVISÃO (em vermelho)
%  \rev{...}   -> trecho ADICIONADO ou CORRIGIDO, em linha.
%               Abre com o adereço  >  e fecha com  <  (ambos vermelhos).
%  ambiente revbox -> observação/parágrafo inteiro em vermelho.
%  Para "aceitar" uma revisão (torná-la texto normal preto), basta
%  trocar \rev{...} pelo conteúdo, ou redefinir \rev para não colorir.
% ---------------------------------------------------------------------
\usepackage{xcolor}
\definecolor{revred}{RGB}{200,0,0}
\newcommand{\rev}[1]{{\color{revred}\ensuremath{\blacktriangleright}\,#1\,\ensuremath{\blacktriangleleft}}}
\newenvironment{revbox}%
  {\par\medskip\noindent\color{revred}\ensuremath{\blacktriangleright}\ \itshape}%
  {\ \upshape\ensuremath{\blacktriangleleft}\par\medskip}

% ---------------------------------------------------------------------
%  ATALHOS DE NOTAÇÃO
% ---------------------------------------------------------------------
\newcommand{\dd}{\mathrm{d}}
\newcommand{\pt}[1]{\partial_{#1}}
\newcommand{\Lag}{\mathcal{L}}
\newcommand{\Ham}{\mathcal{H}}
\renewcommand{\d}{\partial}
\renewcommand{\refname}{Referências}   % título da bibliografia em pt (sem babel)

\title{\textbf{Campos Clássicos}\\[2pt]\large Notas de estudo}
\author{Notas de estudo --- revisadas}
\date{Julho de 2026}

\begin{document}
\maketitle

\begin{revbox}
Nota de revisão. Os trechos em vermelho, entre os adereços
$\blacktriangleright$ e $\blacktriangleleft$, foram \emph{adicionados}
ou \emph{corrigidos} nesta revisão (convenções explicitadas, passagens
completadas, formalismo hamiltoniano, forma geral da corrente de Noether
e referências). O conteúdo em preto reproduz as notas originais,
apenas reorganizado.
\end{revbox}

% =====================================================================
\section{Introdução: o conceito de campo}
% =====================================================================

Vamos começar a analisar \textbf{campos clássicos}, objetos de importância
central para a segunda quantização e para a teoria quântica de campos.

Um \textbf{campo} é um objeto matemático que associa, a cada ponto do
espaço-tempo $x^{\mu}$, uma quantidade --- que pode ser um escalar, um
vetor, um tensor, um espinor, etc. Para um campo escalar,
\begin{equation}
  \phi:\ \mathbb{R}^{1,3}\longrightarrow \mathbb{R}\ (\text{ou } \mathbb{C}),
  \qquad x^{\mu}\longmapsto \phi(x^{\mu}).
\end{equation}

\begin{revbox}
Convenção. Adotamos a métrica de Minkowski
$\eta_{\mu\nu}=\mathrm{diag}(+1,-1,-1,-1)$, $\pt{\mu}\equiv \d/\d x^{\mu}$ e
o operador de d'Alembert $\Box\equiv \pt{\mu}\d^{\mu}=\d_t^{2}-\nabla^{2}$.
Índices repetidos são somados (convenção de Einstein). Fixar a assinatura
é necessário para que os sinais das seções seguintes fiquem consistentes.
\end{revbox}

% =====================================================================
\section{Formalismo lagrangiano}
% =====================================================================

\subsection{Densidade de lagrangiana e ação}

A dinâmica de um campo é descrita pela \textbf{densidade de lagrangiana}
$\Lag(\phi,\pt{\mu}\phi)$, função do campo e de suas derivadas. A
\textbf{ação} é o funcional
\begin{equation}
  S[\phi]=\int_{\Omega}\dd^{4}x\ \Lag(\phi,\pt{\mu}\phi),
  \qquad \dd^{4}x=\dd x^{0}\,\dd x^{1}\,\dd x^{2}\,\dd x^{3}
        =\dd t\,\dd^{3}x .
\end{equation}

\subsection{Equações de Euler--Lagrange}

O \textbf{princípio da mínima ação} exige que a ação seja estacionária,
$\delta S=0$, para variações $\delta\phi$ do campo. Variando,
\begin{equation}
  \delta S=\int_{\Omega}\dd^{4}x
  \left[
    \frac{\d\Lag}{\d\phi}\,\delta\phi
    +\frac{\d\Lag}{\d(\pt{\mu}\phi)}\,\delta(\pt{\mu}\phi)
  \right].
\end{equation}
Como $\delta(\pt{\mu}\phi)=\pt{\mu}(\delta\phi)$, integramos por partes o
segundo termo:
\begin{equation}
  \int_{\Omega}\dd^{4}x\,\frac{\d\Lag}{\d(\pt{\mu}\phi)}\,\pt{\mu}(\delta\phi)
  =\int_{\Omega}\dd^{4}x\,\pt{\mu}\!\!\left[\frac{\d\Lag}{\d(\pt{\mu}\phi)}\,\delta\phi\right]
  -\int_{\Omega}\dd^{4}x\,\pt{\mu}\!\!\left(\frac{\d\Lag}{\d(\pt{\mu}\phi)}\right)\delta\phi .
\end{equation}
O primeiro termo é uma divergência total; pelo teorema de Gauss ele vira
uma integral de superfície na fronteira $\d\Omega$, que se anula
\rev{pois exige-se $\delta\phi\big|_{\d\Omega}=0$ (a variação se anula na
fronteira da região de integração).}
Restando apenas o segundo termo,
\begin{equation}
  \delta S=\int_{\Omega}\dd^{4}x
  \left[\frac{\d\Lag}{\d\phi}
        -\pt{\mu}\!\left(\frac{\d\Lag}{\d(\pt{\mu}\phi)}\right)\right]\delta\phi=0 .
\end{equation}
Como $\delta\phi$ é arbitrário no interior de $\Omega$, o integrando deve
se anular, o que fornece as \textbf{equações de Euler--Lagrange}:
\begin{equation}\boxed{
  \frac{\d\Lag}{\d\phi}
  -\pt{\mu}\!\left(\frac{\d\Lag}{\d(\pt{\mu}\phi)}\right)=0 .}
\end{equation}

% =====================================================================
\section{O campo escalar real: equação de Klein--Gordon}
% =====================================================================

\subsection{Equações de movimento}

Como primeiro exemplo, considere o campo escalar real com densidade de
lagrangiana
\begin{equation}
  \Lag=\tfrac12\,\pt{\mu}\phi\,\d^{\mu}\phi-\tfrac12\,m^{2}\phi^{2}
     =\tfrac12\!\left(\pt{\mu}\phi\,\d^{\mu}\phi-m^{2}\phi^{2}\right).
\end{equation}
O primeiro termo derivado dá
\begin{equation}
  \frac{\d\Lag}{\d\phi}=-m^{2}\phi .
\end{equation}
Para o segundo, escrevemos $\pt{\mu}\phi\,\d^{\mu}\phi
=\eta^{\alpha\beta}\pt{\alpha}\phi\,\pt{\beta}\phi$ e derivamos em relação a
$\pt{\nu}\phi$:
\begin{equation}
  \frac{\d\Lag}{\d(\pt{\nu}\phi)}
  =\tfrac12\,\eta^{\alpha\beta}
   \bigl(\delta^{\nu}_{\alpha}\pt{\beta}\phi+\pt{\alpha}\phi\,\delta^{\nu}_{\beta}\bigr)
  =\tfrac12\bigl(\d^{\nu}\phi+\d^{\nu}\phi\bigr)=\d^{\nu}\phi .
\end{equation}
Substituindo na equação de Euler--Lagrange,
\begin{equation}
  -m^{2}\phi-\pt{\mu}\d^{\mu}\phi=0
  \quad\Longrightarrow\quad
  \boxed{\ (\Box+m^{2})\,\phi=0\ }
\end{equation}
que é a \textbf{equação de Klein--Gordon}.

\subsection{Formalismo hamiltoniano}

Definimos o \textbf{momento canonicamente conjugado} ao campo $\phi$ como
\begin{equation}
  \pi(x)\equiv\frac{\d\Lag}{\d\dot\phi},
  \qquad \dot\phi\equiv\pt{0}\phi,
\end{equation}
e a \textbf{densidade de hamiltoniana} pela transformada de Legendre
\begin{equation}
  \Ham=\pi\,\dot\phi-\Lag .
\end{equation}
Para o campo escalar real, $\pi=\d^{0}\phi=\dot\phi$, e
\begin{revbox}
explicitando a densidade de energia obtemos
\upshape
\begin{equation}
  \color{revred}
  \Ham=\tfrac12\,\pi^{2}+\tfrac12\,(\nabla\phi)^{2}+\tfrac12\,m^{2}\phi^{2},
\end{equation}
\itshape
soma dos termos cinético, de gradiente e de massa --- manifestamente
positiva, como esperado para a energia.
\end{revbox}

% =====================================================================
\section{Teorema de Noether}
% =====================================================================

\subsection{Simetrias e correntes conservadas}

\textbf{Teorema de Noether:} a cada simetria contínua da ação corresponde
uma quantidade conservada. Considere uma transformação infinitesimal do
campo,
\begin{equation}
  \phi(x)\longrightarrow \phi'(x)=\phi(x)+\delta\phi(x),
\end{equation}
que deixa a ação invariante. Variando a lagrangiana e usando as equações
de Euler--Lagrange,
\begin{equation}
  \delta\Lag
  =\frac{\d\Lag}{\d\phi}\,\delta\phi
   +\frac{\d\Lag}{\d(\pt{\mu}\phi)}\,\pt{\mu}(\delta\phi)
  =\pt{\mu}\!\left[\frac{\d\Lag}{\d(\pt{\mu}\phi)}\,\delta\phi\right].
\end{equation}
Se a transformação é uma simetria estrita, $\delta\Lag=0$, e a
\textbf{corrente de Noether}
\begin{equation}
  j^{\mu}=\frac{\d\Lag}{\d(\pt{\mu}\phi)}\,\delta\phi
\end{equation}
é conservada, $\pt{\mu}j^{\mu}=0$.
\begin{revbox}
Caso geral. Se a lagrangiana varia por uma divergência total,
$\delta\Lag=\pt{\mu}K^{\mu}$ (quase-simetria), a ação ainda é invariante e
a corrente conservada passa a ser
\upshape
\begin{equation}
  \color{revred}
  j^{\mu}=\frac{\d\Lag}{\d(\pt{\mu}\phi)}\,\delta\phi-K^{\mu},
  \qquad \pt{\mu}j^{\mu}=0 .
\end{equation}
\itshape
A simetria estrita é o caso particular $K^{\mu}=0$.
\end{revbox}
A cada corrente conservada associa-se uma \textbf{carga conservada}
\begin{equation}
  Q=\int \dd^{3}x\ j^{0},\qquad \frac{\dd Q}{\dd t}=0 .
\end{equation}

\subsection{Tensor de energia--momento}

A invariância sob \textbf{translações no espaço-tempo}
$x^{\mu}\to x^{\mu}+a^{\mu}$ leva à conservação do \textbf{tensor de
energia--momento canônico}
\begin{equation}
  T^{\mu}{}_{\nu}
  =\frac{\d\Lag}{\d(\pt{\mu}\phi)}\,\pt{\nu}\phi-\delta^{\mu}_{\nu}\,\Lag,
  \qquad \pt{\mu}T^{\mu}{}_{\nu}=0 .
\end{equation}
As cargas conservadas correspondentes formam o \textbf{quadrimomento}
\begin{equation}
  P_{\nu}=\int \dd^{3}x\ T^{0}{}_{\nu}.
\end{equation}
Em particular, a componente temporal é a energia total, que reproduz a
hamiltoniana:
\begin{equation}
  H=P^{0}=\int\dd^{3}x\ T^{0}{}_{0}
   =\int\dd^{3}x\ (\pi\dot\phi-\Lag)
   =\int\dd^{3}x\ \Ham .
\end{equation}

% =====================================================================
\section{Simetria \(U(1)\) e o campo escalar complexo}
% =====================================================================

\subsection{Lagrangiana e transformação de fase}

Para um \textbf{campo escalar complexo} $\phi$, tratamos $\phi$ e
$\phi^{*}$ como campos independentes, com densidade de lagrangiana
\begin{equation}
  \Lag=\pt{\mu}\phi^{*}\,\d^{\mu}\phi-m^{2}\phi^{*}\phi .
\end{equation}
\rev{Aqui \emph{não} há o fator $\tfrac12$ presente no caso real: como
$\phi$ e $\phi^{*}$ são graus de liberdade independentes, o fator $\tfrac12$
duplicaria indevidamente os termos. As equações de Euler--Lagrange para
$\phi^{*}$ dão $(\Box+m^{2})\phi=0$ (e a conjugada para $\phi$).}

A lagrangiana é invariante sob a transformação global de fase (grupo
$U(1)$)
\begin{equation}
  \phi\longrightarrow e^{i\alpha}\phi,\qquad
  \phi^{*}\longrightarrow e^{-i\alpha}\phi^{*},
\end{equation}
cuja versão infinitesimal ($\alpha\ll1$) é
\begin{equation}
  \delta\phi=i\alpha\,\phi,\qquad \delta\phi^{*}=-i\alpha\,\phi^{*}.
\end{equation}

\subsection{Corrente conservada e carga}

\rev{Aplicando a corrente de Noether aos dois campos independentes,}
\begin{equation}
  j^{\mu}
  =\frac{\d\Lag}{\d(\pt{\mu}\phi)}\,\delta\phi
   +\frac{\d\Lag}{\d(\pt{\mu}\phi^{*})}\,\delta\phi^{*}
  =(\d^{\mu}\phi^{*})(i\alpha\phi)+(\d^{\mu}\phi)(-i\alpha\phi^{*}).
\end{equation}
Retirando o parâmetro $\alpha$, a corrente conservada é
\begin{equation}
  \boxed{\,j^{\mu}=i\!\left(\phi\,\d^{\mu}\phi^{*}-\phi^{*}\,\d^{\mu}\phi\right),
  \qquad \pt{\mu}j^{\mu}=0.\,}
\end{equation}
\rev{O sinal global é convencional (fixa o sinal da carga); é comum
escrever, de modo equivalente, $j^{\mu}=i(\phi^{*}\d^{\mu}\phi-\phi\,\d^{\mu}\phi^{*})$.}
A carga conservada associada,
\begin{equation}
  Q=\int\dd^{3}x\ j^{0},
\end{equation}
é interpretada, ao se acoplar o campo ao eletromagnetismo, como a
\textbf{carga elétrica}.

% =====================================================================
\begin{revbox}
Seção de referências adicionada na revisão (formato ABNT NBR 6023).
\end{revbox}
\begin{thebibliography}{9}
\bibitem{peskin}
  PESKIN, M. E.; SCHROEDER, D. V.
  \textbf{An introduction to quantum field theory}.
  Reading: Addison-Wesley, 1995.

\bibitem{weinberg}
  WEINBERG, S.
  \textbf{The quantum theory of fields}: foundations. v.~1.
  Cambridge: Cambridge University Press, 1995.

\bibitem{greiner}
  GREINER, W.; REINHARDT, J.
  \textbf{Field quantization}.
  Berlin: Springer, 1996.

\bibitem{goldstein}
  GOLDSTEIN, H.; POOLE, C.; SAFKO, J.
  \textbf{Classical mechanics}. 3.~ed.
  San Francisco: Addison-Wesley, 2002.
\end{thebibliography}

\end{document}
